\section{Introduction}
\label{sec:introduction}

For a finite-dimensional algebra \(A\) over a field \(k\), its
\emph{little finitistic dimension} is
\[
\findim A
=\sup\{\pd_A N:
       N\text{ is a finitely generated left }A\text{-module},
       \ \pd_A N<\infty\}.
\]
The little finitistic-dimension conjecture asserts that this number
is finite for every such algebra.  Thus it asks for a bound on the
lengths of the finite projective resolutions over each fixed algebra,
while allowing other modules to have infinite projective dimension.
The finitistic-dimension questions were publicized by Bass in
1960 \cite{Bass1960}.  The \emph{big} finitistic dimension \(\Findim A\)
is defined by allowing all left modules in the same supremum; the little
dimension is at most the big one.

This paper gives a negative resolution of the little
finitistic-dimension conjecture.

\begin{theorem}\label{thm:main}
There is a finite-dimensional unital complex algebra \(A\) such that,
for every integer \(m\ge1\), there is a finite-dimensional left
\(A\)-module \(N_m\) with
\[
2m-2\le\pd_A N_m<\infty.
\]
In particular, \(\findim A=\infty\).
\end{theorem}

The algebra in \Cref{thm:main} is fixed before \(m\) varies.
Consequently its finitely generated modules admit terminating
projective resolutions of unbounded lengths.  This also makes its
big finitistic dimension infinite and refutes the corresponding
universal finiteness assertion for Artin algebras, since a
finite-dimensional complex algebra is an Artin algebra.

Independently, the companion~\cite[Corollary~1.3]{OpenAITachikawa2026}
obtains finite-dimensional opposite endomorphism algebras \(\Gamma_K\)
of infinite little and big left finitistic dimensions for every field
extension \(K/k\), where \(k=\mathbb F_2(q,H_1,H_2)\), using successive
cokernels in a projective-injective resolution. This is a distinct
characteristic-two example; the algebra in Theorem~\ref{thm:main}
remains the complex example constructed here.

\subsection{Background and consequences}

The equality of the little and big finitistic dimensions is a different
question from their finiteness. Huisgen-Zimmermann disproved that equality:
for every integer \(n\ge2\), she constructed a finite-dimensional monomial
algebra with little finitistic dimension \(n\) and big finitistic dimension
\(n+1\) \cite{HuisgenZimmermann1992}. Those examples have finite finitistic
dimensions and vary with \(n\); they do not give unbounded finite projective
dimensions over a single algebra. Her survey \cite{HuisgenZimmermann1995}
describes the development of the two questions.

Important positive cases for little finitistic finiteness include
finite-dimensional monomial algebras,
by Green, Kirkman, and Kuzmanovich \cite{GreenKirkmanKuzmanovich1991},
and Artin algebras with radical cube zero, by Green and
Huisgen-Zimmermann \cite{GreenHuisgenZimmermann1991}.
Igusa and Todorov introduced numerical invariants of syzygies that give
another proof of the radical-cube-zero case and establish little
finiteness for Artin algebras of representation dimension at most three
\cite[Corollaries~7 and~9]{IgusaTodorov2005}.
More recently, Cummings proved that universal little finitistic
finiteness is equivalent to the assertion that, for every
finite-dimensional algebra, finiteness holds on the left if and only
if it holds on the right \cite[Theorem~A]{Cummings2024}.
These results address bounds within classes or reductions between
finiteness assertions. Our construction instead preserves arbitrarily
long terminating processes while keeping all defining algebra data fixed.

There is also a derived-category consequence. Say that injective left
modules \emph{generate} if the smallest triangulated subcategory of the
unbounded derived category of all left modules that contains them and is
closed under arbitrary coproducts is the whole category. Rickard proved
that injective generation implies finite big finitistic dimension
\cite[Theorem~4.3]{Rickard2019}. Thus injective left \(A\)-modules do not
generate for the algebra of \Cref{thm:main}: its big left finitistic
dimension is infinite. This applies Rickard's theorem to
\(A^{\mathrm{op}}\), matching his right-module convention.

Combining this example with Cummings' auxiliary triangular algebra gives
an extreme left--right asymmetry for both dimensions and for injective
generation.

\begin{corollary}[Extreme left--right asymmetry]
\label{cor:left-right-asymmetry}
There is a finite-dimensional unital complex algebra \(\Lambda\) such that
\[
\findim \Lambda=\Findim \Lambda=\infty,
\qquad
\findim(\Lambda^{\mathrm{op}})=\Findim(\Lambda^{\mathrm{op}})=0.
\]
All four dimensions here use left modules. Thus the little and big
finitistic dimensions of \(\Lambda\) are infinite on the left and zero
on the right. Injective left \(\Lambda\)-modules do not generate the
unbounded derived category of all left \(\Lambda\)-modules, whereas
injective right \(\Lambda\)-modules generate the unbounded derived
category of all right \(\Lambda\)-modules.
\end{corollary}

The algebra \(\Lambda\) is auxiliary: the corollary makes no assertion
about the right finitistic dimensions or right injective generation of
the original algebra \(A\). Passing from \(\Lambda\) to
\(\Lambda^{\mathrm{op}}\) reverses the two sides of the corollary.
We give the proof in \Cref{sec:left-right-asymmetry}, after constructing
\(A\).

\subsection{The construction and its principal mechanisms}

The proof connects a finite presentation with a derived tensor
process and then with ordinary projective resolutions.
We describe the three stages, including the points at which
compatibility of the constructions is needed.

First, \Cref{prop:selection} constructs one finitely presented
complex algebra \(R\), a central idempotent \(e\), and an automorphism
\(\alpha\).  The functor
\[
H(Y)={}_{\alpha}(eY)
\]
selects the direct summand \(eY\) and restricts its action along
\(\alpha\).  It admits finite-dimensional modules \(Y_m\) that
survive exactly \(m-1\) iterations and vanish on the \(m\)th.
The algebra \(R\) is a group algebra.  In the underlying group,
an automorphism shifts an indexed family of central involutions;
finite matrix quotients detect arbitrarily long independent
segments of that family.  Prescribing their characters gives the
modules \(Y_m\).  The group is motivated by matrix groups of
Abels type \cite{Rego2022}; its finite presentation, centrality
relations, quotients, and representations are established here.

The second stage realizes this selection process over a
finite-dimensional algebra.  A quadratic presentation of \(R\)
gives a directed algebra \(B\) with three vertices and
\(\gldim B\le2\), together with modules \(M(Y)\).
The module \(M(Y)\) places the vector space \(Y\) at each vertex and
records the generator actions along the arrows.
The encoding follows the homogenization principle used in
universal localization \cite{NRS2004}.
The principal realization statement, \Cref{thm:realization},
constructs a fixed bounded \(B\)-bimodule complex \(P\), projective
term by term on the right, for which
\[
P\otL_B M(Y)\simeq M(H(Y))\oplus M(H(Y))[3]
\]
for every finite-dimensional \(R\)-module \(Y\).

The principal obstacle is to make this realization uniform in \(Y\).
Inverting the two distinguished arrows of \(B\) makes the presentation
of \(R\) act on a localized category of complexes. The summand selected
by \(e\) need not lift by itself; two mapping cones represent its
direct sum with a copy shifted by three inside that localized category
(\Cref{lem:odd-double}).
We then clear denominators for one finite diagram, before choosing any
test module (\Cref{lem:diagram-lift}). Its relations initially hold only
up to chain homotopy. A three-column differential incorporates those
homotopies into an actual bimodule complex (\Cref{lem:rectification}).
There are no paths of length three, so no further relation overlaps
require higher compatibility data. Finally, the gap between the two
surviving cohomological degrees exceeds the global-dimension bound on
\(B\), so the complex splits as displayed. These steps realize selection
by a single tensor functor over \(B\).

The third stage turns this bounded complex into an ordinary
bimodule.  A directed radical-square-zero chain records the terms
and differential of \(P\).  With a suitable product algebra \(D\)
of finite global dimension and an ordinary \(D\)-bimodule \(X\),
the second iterate of \(F=X\otL_D-\) is isomorphic, on the
\(B\)-factor of \(D\), to \(P\otL_B-\) up to a fixed shift
(\Cref{prop:ordinary-simulation}).  We then take
\[
A=D\ltimes X.
\]
For a \(D\)-module \(N\) inflated to \(A\), the absolute bar
construction decomposes by the number of occurrences of \(X\):
\[
D\otL_A N\simeq\bigoplus_{r\ge0}(F^rN)[r].
\]
\Cref{prop:bar-splitting} proves this decomposition with derived
tensor powers, including the nonflat case. This splitting is also a
consequence of Minamoto and Yamaura's graded-component decomposition
for trivial extensions
\cite[Lemma~4.13(4) and the proof of Theorem~4.17]{MinamotoYamaura}.
Their projective-dimension formula gives a stronger form of the final
transfer. We prove the required splitting directly by the bar construction
and give the minimal-resolution argument used here.
Eventual vanishing of \(F^rN\) yields finite projective dimension,
while a nonzero \(r\)th iterate forces projective dimension at
least \(r\). For \(N_m=M(Y_m)\), the two-step simulation of selection
gives \(F^{2m}N_m\simeq0\) and \(F^{2m-2}N_m\not\simeq0\).
These yield the finiteness and lower bound in \Cref{thm:main}, respectively.

\Cref{thm:realization} applies to any finitely presented
\(k\)-algebra with a central idempotent and a unital endomorphism.
Readers interested in the homological constructions may begin with
\Cref{sec:localization} and use \Cref{prop:selection} as input;
\Cref{sec:selection} supplies the unbounded extinction times needed
for the final counterexample.

In the proof order, \Cref{sec:selection} supplies the selection modules;
\Cref{sec:localization} constructs the fixed lifted diagram; and
\Cref{sec:realization} produces and identifies its bimodule complex.
\Cref{sec:ordinary-simulation} makes the bimodule ordinary, and
\Cref{sec:square-zero} turns its iterates into projective-dimension bounds.

\subsection{Conventions}

All algebras are associative and unital.  Modules are left modules
unless expressly designated right; over a finite-dimensional
\(k\)-algebra, finite dimensionality and finite generation are
equivalent.  Tensor products carrying a superscript \(\mathbf L\)
are derived tensor products.  Complexes have cohomological grading,
with
\[
K[s]^n=K^{n+s},
\qquad d_{K[s]}^n=(-1)^s d_K^{n+s}.
\]
In particular, a module shifted by \(s\) has its cohomology in
degree \(-s\).  Products of paths and endomorphisms are written
in order of composition, from right to left.
An isomorphism denoted \(\simeq\) is taken in the derived category
under discussion.  Splittings used to iterate a tensor functor
are needed only for individual objects; no natural choice of
those splittings is asserted.
